Small language models are usually judged by how much they know.
We ask how little a model needs to know if the rest can be read from its context.
We frame this as the search for a \emph{base curriculum}: the minimal set of skills that must be in the weights so that task-specific knowledge can be supplied as instructional material at inference time.
We propose a six-layer candidate, a criterion for what belongs in the weights, and a ladder of generalisation that separates memorising a textbook from reading one.
We introduce TinyTextbook, a synthetic benchmark in which each episode is a small textbook of freshly sampled rules followed by exercises, and report a pilot study with transformers of about \NparamsRounded{} parameters trained on a CPU (\TotalRuns{} runs).
Models trained on 16 or fewer distinct textbooks answer perfectly on those and at or near chance on new ones; with enough distinct textbooks they read new ones.
A skill left out of training stays at chance, and without the simplest skill the models did not acquire composition either.
Separately trained skills did not combine: no run exceeded \ETwoMixedMax\% on exercises mixing two of them (guessing: 26.7\%).
Applying the same two layers twice extended the chain of rules a model could follow without adding parameters, but at equal depth looped and unshared models showed no consistent difference.
Outcomes vary between seeds and are reported as counts.
The pilot is small: one to four seeds, one synthetic domain, and a fixed page layout in which values are found by position; retrieval by content was not learned within our budget.
We give a protocol with falsifiable hypotheses for 10--100 million parameters and release the code.
